\begin{textsample}{POS Dim 3 – global\_north\_2024\_03 – Score 10.00 – global\_north\_2024\_03/200310789.txt}  \label{ex:f3_pos_276}
How one NJ \textbf{college} is fostering dialogue over Israeli-Palestinian conflict Hannan Adely, NorthJersey.com Updated March 18, 2024 at 1:58 PM*3 min read As tensions flare on \textbf{college} \textbf{campuses} over the Israeli-Palestinian conflict, one North Jersey \textbf{school} is trying to bridge the divide with a series of events aimed at humanizing"the other"and promoting dialogue.
Ramapo \textbf{College} ' s program, called"Power of Conversation,"is bringing together experts, \textbf{students} and individuals with firsthand reports from the region.
The \textbf{college} in Mahwah hosted speakers who delivered accounts of the Oct. 7 attack in Israel and of the dire situation and needs on the ground in Gaza.
On Tuesday, in the third talk in the series, two scholars gave insights on how to approach dialogue on a subject that can feel fraught with emotion to some and intimidating for others, who fear getting it wrong or offending people.
Here are three takeaways from the presentation, called"Meaningfully Discussing Palestine/Israel after October 7,"and interviews with Mira Sucharov, a professor of political science at @ @ @ @ @ @ @ @ @ @ of peace and conflict studies and anthropology at Swarthmore \textbf{College}.
Get educated Look at a wide range of sources from American and international press to get an understanding of current events.
Look at a wide range of sources from the American and international press to get an understanding of current events, Atshan advised.
He reads the publication Haaretz daily, he said, adding that Israeli media give more voice to Israel ' s critics than American mainstream news organizations so."It ' s useful to, alongside CNN, to look at the Israeli press, to look at the Arab press, to look at Al Jazeera... so it ' s curating a combination of social media sources as well as mainstream media sources with a wide range of points of view and wide range of regions,"Atshan said.
Sucharov recommended explainers on the conflict available on the Vox website to get a grasp of the history and politics of the region.
Humanize the other side Two customers share a meal at BurgerHaus @ @ @ @ @ @ @ @ @ @"Sucharov said.
The Canadian professor and author encouraged people to share their concerns and fears, then say what they need to feel safe.
The conversation can start with areas of agreement and recognizing one another ' s pain."Even if you find you are clashing a lot, you say, well, let ' s start a new round of conversation today and let ' s start on what we can both agree on, and that builds a common ground,"she said.
She also suggested that people or groups having tough conversations also spend time in other ways, such as sharing meals or playing sports or"something that just humanizes you and gets you to see the other as three-dimensional human beings."Those relationships, she said,"can just reduce the tension and create trust."Protect free \textbf{speech} A flier posted inside the Jon Huntsman Hall at the \textbf{University} of Penn ' s Wharton \textbf{School} takes the \textbf{university} to task for its stance on free \textbf{speech}.
Colleges and @ @ @ @ @ @ @ @ @ @ paramount, even as they face intense external pressures to quell \textbf{speech} tied to the Israeli-Palestinian conflict, both speakers said Tuesday."When \textbf{universities} are getting donor dollars, they have to be very clear that that donor realizes that \textbf{academic} freedom is part of the deal,"Sucharov said.
Atshan said \textbf{institutions} have to stand behind their values, principles and commitments to free \textbf{speech}."\textbf{Faculty} in the classroom and \textbf{students} in the classroom should feel that they can engage in ideas, engage in critical thinking and play devil ' s advocate without negative repercussions, especially from external groups that are pressuring \textbf{campuses} in a McCarthyist way to try to \textbf{chill} free \textbf{speech},"he said."I think we have to stand up to those bullies."What ' s next?
Ramapo \textbf{College} held a \textbf{student} forum Thursday called"I Think, I Feel,"described as a space for \textbf{students} to share their thoughts about current issues in the United States and the world and listen to different viewpoints. @ @ @ @ @ @ @ @ @ @ installation called"Be Heard.
Listen."Viewers can select icons or photos to hear \textbf{students} ' answers to prompts and questions that shed light on their experiences and perspectives.
Recommended Stories The Hamas-controlled Gaza health ministry said Israeli forces opened fire on the crowd.
Israeli military officials said its troops fired warning shots before firing"only in face of danger when the mob moved in a manner which endangered them."Users of Elon Musk-owned X ( formerly Twitter ) continue complaining the platform is engaging in shadowbanning -- aka restricting the visibility of posts by applying a"temporary"label to accounts that can limit the reach/visibility of content -- without providing clarity over why it ' s imposed the sanctions.
Many users can be seen expressing confusion as to why they ' re being penalized -- apparently not having been given a meaningful explanation as to why the platform has imposed restrictions on their content.
Complaints that surface in a search for the phrase"temporary label"show users appear to have received only generic notifications @ @ @ @ @ @ @ @ @ @ text in which X states their accounts"may contain spam or be engaging in other types of platform manipulation".

% matched lemmas: academic, campus, chill, college, faculty, institution, school, speech, student, university
\end{textsample}
